% QMB_n07_chsh_experiment.tex — Kapitel 7: ΔCHSH und Experiment
% Quellen: Dok. 190, 230, 147
\chapter{CHSH-Abweichung und experimentelle Reichweite}
\label{ch:chsh_exp}

\section{Zwei verschiedene Observablen}

Das Korrekturregister (Dok.~190, Präzisierung~11 und 12, Mai 2026)
unterscheidet präzise zwischen zwei verschiedenen Größen, die frühere
Dokumente zusammenwarfen:

\textbf{Elementare Abweichung.} Pro einzelner Messung entstehen
Phasenkosten des $\theta_4$-Projektionszyklus. Der natürliche
Kandidat ist:
\begin{equation}
  \Delta\mathrm{CHSH}_{\mathrm{elem}} = \frac{\xi}{2\pi}
  \approx 2{,}1\times10^{-5}.
\end{equation}
Diese Identifikation ist eine Arbeitshypothese; eine formale Herleitung
aus der $T^4$-Projektionsgeometrie fehlt noch.\status{S}

\textbf{Kumulative Abweichung.} Über $N$ Messläufe wächst die
Abweichung nach der Formel aus Dok.~022:
\begin{equation}
  \Delta\mathrm{CHSH}(N) = 2\sqrt2 \Bigl[1 - \exp\!\Bigl(-\xi\,\frac{\ln N}{D_f}\Bigr)\Bigr]
  \approx \xi\,\frac{\ln N}{D_f}\cdot2\sqrt2.
\end{equation}
Bei $N=73$ ergibt das $\approx5\times10^{-4}$. Diese Größe ist
experimentell zugänglich, sobald Fehlerkorrektur-Schwellenwerte
unterschritten sind.

\section{Experimentelle Situation 2026}

IBM Heron~r2 (Mai 2026, Dok.~147) liefert:
\begin{itemize}
  \item Bell-Verletzung bestätigt: $S=2{,}74$ ($96{,}9\,\%$ der
        Tsirelson-Schranke).\status{K}
  \item FFGFT-Schaltkreise lauffähig auf aktueller Hardware.\status{K}
  \item Geräte-zu-Geräte-Variation: $1{,}4\,\%$ — rund
        $1400\times$ größer als der gesuchte $\xi$-Effekt.\status{K}
\end{itemize}

Der $\xi$-Effekt erfordert sub-ppm-stabile fehlerkorrigierte Qubits.
Mit NISQ-Hardware nicht auflösbar.\status{K}

\section{Offene Herleitung}

Die Identifikation $\Delta\mathrm{CHSH}_{\mathrm{elem}} = \xi/(2\pi)$
ist plausibel, aber nicht aus der $T^4$-Projektionsgeometrie hergeleitet.
Das Verhältnis zwischen kumulativer und elementarer Abweichung bei
$N$ Läufen beträgt:
\begin{equation}
  \frac{\Delta\mathrm{CHSH}(N)}{\Delta\mathrm{CHSH}_{\mathrm{elem}}}
  \approx N\cdot\frac{2\pi}{D_f},
\end{equation}
bei $N=73$ einen Faktor $\approx153$, konsistent mit dem beobachteten
Faktor $\sim50$--$100$.\status{S}

Die formale Herleitung von $\xi/(2\pi)$ aus der Trägergeometrie
bleibt eine offene Brücke.

\quelle{Dok.~022, 147, 190 (Präz.~11, 12), 230}
